\section*{Capítulo \ref{ch:b2:waves-in-media} --- Ondas eletromagnéticas em plasmas, condutores e dielétricos}

\begin{solution}{exo:b2:waves-in-media:1}
$f_p \approx 9\sqrt n$ Hz: \qty{2.8}{MHz}, \qty{9}{MHz}, \qty{9}{GHz}, \qty{2.6e15}{Hz}
($\lambda_p = \qty{115}{nm}$), \qty{2.8}{kHz}. AM a \qty{1}{MHz}: só pelo
espaço interestelar. FM: pela ionosfera, mas não pela bainha. GPS:
por tudo, menos pela bainha. Luz: por tudo, menos pelo cobre.
\end{solution}

\begin{solution}{exo:b2:waves-in-media:2}
Cobre: \qty{9.2}{mm}, \qty{2.1}{mm}, \qty{65}{\micro m}, \qty{2.1}{\micro m}; alumínio a
\qty{1}{MHz}: \qty{84}{\micro m}; água do mar: \qty{2.3}{m} a \qty{10}{kHz}, \qty{7}{mm} a
\qty{1}{GHz}; solo: \qty{5}{m} a \qty{1}{MHz}. Um submarino a \qty{10}{m} precisa de $\delta
\gtrsim \qty{10}{m}$: abaixo de cerca de \qty{500}{Hz}.
\end{solution}

\begin{solution}{exo:b2:waves-in-media:3}
$n = 1.528$, $1.515$, $1.509$; $v = 1.963$, $1.980$, $\qty{1.987e8}{m/s}$; $n_g = n +
2B/\lambda^2 = 1.545$. Ângulos de refração $\arcsin(0.707/n)$: \qty{27.57}{\degree} e
\qty{27.93}{\degree}: uma abertura de \qty{0.36}{\degree}.
\end{solution}

\begin{solution}{exo:b2:waves-in-media:4}
$f_p = \qty{4.0}{MHz}$: \qty{5}{MHz} passa; \qty{2}{MHz} é refletida, penetrando
$c/\sqrt{\omega_p^2 - \omega^2} = \qty{14}{m}$. Para $n = \qty{1.5e12}{m^{-3}}$: \qty{11}{MHz}. A
densidade acompanha o fluxo ionizante do Sol --- hora, estação e o
ciclo de onze anos ---, de modo que a maior frequência utilizável se desloca com eles.
\end{solution}

\begin{solution}{exo:b2:waves-in-media:5}
(a) $1/\gamma\pi a^2 = \qty{5.3}{m\ohm/m}$. (b) $\delta = \qty{65}{\micro m}$; $R \approx 1/\gamma2\pi a\delta
= \qty{41}{m\ohm/m}$, $7.7$ vezes a de contínua. (c) $\delta = \qty{6.5}{\micro m}$, \qty{0.41}{\ohm/m},
$77$ vezes. (d) $R \propto \sqrt f$ cresce tão rápido quanto $L\omega$ quando o \omterm{thm:b2:waves-in-media:skin}{efeito
pelicular} domina, e cada espira acrescenta resistência; o fio Litz reúne muitos
filamentos isolados mais finos que $\delta$, de modo que todo o cobre conduz.
\end{solution}

\begin{solution}{exo:b2:waves-in-media:6}
$\delta_{\text{Al}} = 11.9$, $0.84$, \qty{0.084}{mm}: $d/\delta = 0.08$, $1.2$, $12$: $0.7$,
$10$, \qty{103}{dB}. A \qty{50}{Hz} o condutor é transparente; o ferro de alto
$\mu_r$ desvia o fluxo (blindagem magnética) e ainda encolhe $\delta$ por
$\sqrt{\mu_r}$.
\end{solution}

\begin{solution}{exo:b2:waves-in-media:7}
(a) $v_0 = eE_0/m\omega$; $\tfrac12nm\langle v^2\rangle = ne^2E_0^2/4m\omega^2 = \tfrac14\varepsilon_0
E_0^2\,\omega_p^2/\omega^2$. (b) $\tfrac14\varepsilon_0E_0^2$ e $\tfrac14\varepsilon_0E_0^2(1 - \omega_p^2/
\omega^2)$. (c) A soma vale $\tfrac12\varepsilon_0E_0^2$ (o parêntese vale $1$); $\langle\Pi
\rangle = E_0^2k/2\mu_0\omega$, e a razão dá $c^2k/\omega = v_g$. (d) $\tfrac14\cdot\tfrac14
/\tfrac12 = 1/8$.
\end{solution}

\begin{solution}{exo:b2:waves-in-media:8}
(a) $1/v_g = (1/c)(1 - \omega_p^2/\omega^2)^{-1/2} \approx (1 + \omega_p^2/2\omega^2)/c$. (b) $\Delta t
= (L\omega_p^2/2c)(1/\omega_1^2 - 1/\omega_2^2)$ com $\omega_p^2 = ne^2/\varepsilon_0m$ e $\omega =
2\pi f$: constante $e^2/8\pi^2\varepsilon_0mc = \qty{1.34e-7}{}$ (SI). (c) $nL = \qty{9.3e23}{m^{-2}}$,
$(1/f_1^2 - 1/f_2^2) = \qty{5.7e-18}{s^2}$: $\Delta t = \qty{0.72}{s}$. (d) A coluna $nL$
(a medida de dispersão) e, com um modelo de $n$, a distância.
\end{solution}

\begin{solution}{exo:b2:waves-in-media:9}
(a) $\operatorname{Im}\underline\chi = \omega_p^2\Gamma\omega/[(\omega_0^2 - \omega^2)^2 + \Gamma^2\omega^2]$ com
$\omega_0^2 - \omega^2 \approx 2\omega_0(\omega_0 - \omega)$: $n'' = \tfrac12\operatorname{Im}\underline\chi$ é
a lorentziana enunciada. (b) $n''(\omega_0) = \omega_p^2/2\omega_0\Gamma$, $\alpha = 2n''\omega_0/c =
\omega_p^2/\Gamma c$. (c) $\omega_p^2 = \qty{3.2e20}{s^{-2}}$: $\alpha = \qty{1.8e4}{m^{-1}}$;
$\ln10/\alpha = \qty{0.13}{mm}$. (d) $n' - 1 \propto (\omega_0 - \omega)/[(\omega_0 - \omega)^2 +
\Gamma^2/4]$: positivo abaixo, negativo acima, decrescente dentro de $\pm\Gamma/2$
da linha --- anômala ali.
\end{solution}

\begin{solution}{exo:b2:waves-in-media:10}
(a) $\underline E = E_0\eu^{-x/\delta}\eu^{\iu(\omega t - x/\delta)}$, $\underline B = \underline k\underline E/\omega
= (1 - \iu)\underline E/\omega\delta$. (b) $\langle\Pi_x\rangle = \tfrac12\operatorname{Re}(\underline E
\underline B^*)/\mu_0 = E_0^2/2\mu_0\omega\delta$ em $x = 0$. (c) $\int_0^\infty\tfrac12\gamma E_0^2
\eu^{-2x/\delta}\dd x = \gamma E_0^2\delta/4 = E_0^2/2\mu_0\omega\delta$, usando $\gamma = 2/\mu_0\omega\delta^2$.
(d) $H_0 = \sqrt2E_0/\mu_0\omega\delta$, de modo que $\tfrac12R_sH_0^2 = E_0^2/\gamma\mu_0^2\omega^2\delta^3 =
E_0^2/2\mu_0\omega\delta$. Cobre a \qty{1}{GHz}: $R_s = 1/\gamma\delta = \qty{8}{m\ohm}$;
$\tfrac12 \times 8 \times 10^{-3} \times 10^4 = \qty{40}{W/m^2}$.
\end{solution}

\begin{solution}{exo:b2:waves-in-media:11}
(a) $\omega_p = \qty{1.4e16}{rad/s}$, $\lambda_p = 2\pi c/\omega_p = \qty{140}{nm}$. (b) $\omega = \qty{3.8e15}{rad/s}
< \omega_p$: $1/\kappa = c/\sqrt{\omega_p^2 - \omega^2} = \qty{23}{nm}$, reflexão total. (c)
Ainda abaixo de $\omega_p$: refletora; a transparência só além de \qty{140}{nm} (no
sódio, \qty{210}{nm}). (d) As colisões dão uma parte real a $\underline\gamma$ e
alguns por cento de absorção; as ressonâncias de elétrons ligados no visível
(ouro, cobre) absorvem o azul e colorem o metal.
\end{solution}

\begin{solution}{exo:b2:waves-in-media:12}
(a) $\omega\tau = 0.12$: $\underline\varepsilon_r = 84 - 9.7\iu$; a \qty{20}{GHz}, $\omega\tau = 1$:
$45 - 40\iu$. (b) $\underline n \approx 9.2 - 0.53\iu$; $c/2n''\omega = \qty{1.8}{cm}$: o forno
cozinha os dois centímetros externos, e a condução faz o resto. (c) $\tfrac12 \times
1.54 \times 10^{10} \times 8.85 \times 10^{-12} \times 9.7 \times 4 \times 10^6 = \qty{2.6}{MW/m^3}$. (d) A
\qty{20}{GHz} o comprimento de absorção seria de um milímetro: só a casca
cozinharia; \qty{2.45}{GHz} penetra e é uma faixa livre.
\end{solution}

\begin{solution}{pb:b2:waves-in-media:1}
\textbf{1.} $f_p = \qty{9}{MHz}$: o GPS passa, a onda curta é refletida.

\textbf{2.} $k = (\omega/c)\sqrt{1 - \omega_p^2/\omega^2}$: $v_\varphi = \omega/k \approx c(1 + \omega_p^2/
2\omega^2)$, $v_g = \dd\omega/\dd k = c^2k/\omega \approx c(1 - \omega_p^2/2\omega^2)$.

\textbf{3.} $\Delta t = \int(1/v_g - 1/c)\dd z = \int\omega_p^2\dd z/2c\omega^2 = e^2N_T/2\varepsilon_0
mc\omega^2 = (e^2/8\pi^2\varepsilon_0mc)N_T/f^2 = 1.34 \times 10^{-7} \times 2 \times 10^{17}/2.48 \times
10^{18} = \qty{11}{ns}$: \qty{3.2}{m}.

\textbf{4.} $N_T = (\Delta t_1 - \Delta t_2)/K(1/f_1^2 - 1/f_2^2)$; a diferença vale
\qty{7}{ns}: para $1\%$ é preciso \qty{0.07}{ns}.

\textbf{5.} $v_\varphi$ excede $c$ tanto quanto $v_g$ fica aquém: a
fase da portadora chega adiantada, enquanto a modulação chega atrasada. Um
receptor de fase de portadora aplica a correção com o sinal oposto.

\textbf{6.} \qty{3.2}{m} de dia, \qty{0.3}{m} de noite.

\textbf{7.} $f_p = \qty{0.9}{MHz}$ com colisões: as ondas médias e curtas são
absorvidas (um apagão de rádio); o GPS a \qty{1.5}{GHz} sofre uma absorção
$\propto 1/\omega^2$, desprezível.

\textbf{8.} $f_1\Delta t = 1.575 \times 10^9 \times 1.1 \times 10^{-8} = 17$ ciclos.

\textbf{9.} $\delta = \sqrt{2/\mu_0\gamma\omega} = \qty{10}{\micro m}$: uma parede de um milímetro tem
cem \omterm{def:b2:dispersion-wave-packets:complex}{profundidades de penetração}.

\textbf{10.} $R_s = 1/\gamma\delta = \qty{0.1}{\ohm}$; $\tfrac12R_sH_0^2 = \qty{45}{W/m^2}$: \qty{45}{W},
$6\%$.

\textbf{11.} $\eu^{-\pi \times 0.5} = 0.21$: \qty{-14}{dB} em amplitude e \qty{-27}{dB} em
potência (e o furo minúsculo já acopla pouco, para começar).

\textbf{12.} A luz, $\lambda = \qty{0.5}{\micro m}$, é duas mil vezes menor
que os furos e passa; a micro-onda, $\lambda = \qty{12}{cm}$, é
cem vezes maior.

\textbf{13.} $\underline n = 2 - 0.005\iu$; $c/2n''\omega = \qty{1.9}{m}$: o vidro quase
não é aquecido.

\textbf{14.} $j \approx H_0/\delta = \qty{3e7}{A/m^2}$, $j^2/\gamma = \qty{9e8}{W/m^3}$ numa
camada de dez micrômetros: as pontas aquecem em milissegundos, emitem elétrons,
ionizam o ar --- faíscas.

\textbf{15.} $\delta = \qty{18}{\micro m}$; seção $2\pi a\delta = \qty{5.6e-8}{m^2}$ contra
$\pi a^2 = \qty{7.9e-7}{m^2}$: $R_{\text{ca}} = \qty{0.30}{\ohm/m}$, catorze vezes a de contínua.

\textbf{16.} $RI^2 = \qty{120}{W/m}$; superfície $\pi d = \qty{3.1e-3}{m^2}$ por metro:
$\Delta T \approx \qty{1900}{K}$ --- ele fundiria; essas bobinas são tubos refrigerados a água.

\textbf{17.} $\delta = \qty{0.14}{mm}$: a corrente induzida, e o calor, ficam
num décimo de milímetro; um pulso curto seguido de têmpera endurece
a superfície de uma engrenagem, enquanto o núcleo permanece tenaz.

\textbf{18.} $R = 1/\gamma2\pi a\delta = \qty{0.10}{\ohm/m}$, três vezes menos; o
interior não conduz nada --- daí um tubo, com o líquido de refrigeração onde o
cobre seria desperdiçado.

\textbf{19.} $\delta = a$ em $f = 1/\pi\mu_0\gamma a^2 = \qty{17}{Hz}$; abaixo disso a corrente é
uniforme e vale a fórmula de contínua.

\textbf{20.} $n^2 \approx 1 + \omega_p^2/\omega_0^2 + (\omega_p^2/\omega_0^2)(\lambda_0/\lambda)^2$: a forma de Cauchy
com $n(\infty)^2 - 1 = \omega_p^2/\omega_0^2 = 1.25$ para $n(\infty) = 1.5$.

\textbf{21.} $n = 1.531$ e $1.510$, $v = 1.96$ e \qty{1.99e8}{m/s}; $n_g(550) =
1.550$; $n_g(400) - n_g(700) = 0.063$: \qty{210}{ns} por quilômetro.

\textbf{22.} $\sin((A + D)/2) = 1.5165 \times 0.5$: $D = \qty{38.6}{\degree}$; $\dd D = 2\sin
(A/2)\dd n/\cos((A + D)/2) = 1.53\,\dd n$: $\Delta D = \qty{1.85}{\degree}$, \qty{6.5}{cm} a \qty{2}{m}.

\textbf{23.} A \qty{200}{nm} a frequência se aproxima da ressonância e
$n''$ cresce: opaco. Uma ressonância fraca a \qty{700}{nm} absorve o vermelho:
vidro verde.

\textbf{24.} O azul está mais perto da ressonância do ultravioleta: os elétrons
ligados respondem mais, $n$ é maior, e o desvio também.

\textbf{25.} Camada F a \qty{1.5}{GHz}: $\underline\gamma$ imaginária, propagação
com pequeno atraso de grupo. Aço a \qty{2.45}{GHz}: $\underline\gamma$ real, \omterm{thm:b2:waves-in-media:skin}{efeito
pelicular}, reflexão. Cobre a \qty{13.56}{MHz}: o mesmo, com uma película de
\qty{18}{\micro m}. Vidro no visível: $\iu\omega\varepsilon_0\chi$, índice real, um
\omterm{def:b2:dispersion-wave-packets:relation}{meio dispersivo} transparente.
\end{solution}
